% Medorama's HTTP API, version 1: the conventions (sign-in, the change header, answers, errors,
% Privacy Mode) and one reference entry per route. The authority is medorama/src/server.rs (the
% router at its end) and the handlers it calls; medorama/openapi.yaml describes the same routes
% for tools, and a test keeps it equal to the router.
\chapter{The HTTP API}
\label{chap:http-api}

\md{}'s page and every other client talk to the server through one JSON API, version 1, under
\file{/api/v1/}. This chapter is for integrators who write scripts or tools against that API. It
describes how to sign in, the rules that apply to every request, and every route, one entry per
route with its parameters, its answer and its errors. The same routes are described for tools in
\file{medorama/openapi.yaml} (OpenAPI 3.1); a test of the repository fails when a route is added,
removed or changed on one side only.

In the entries, a path is written relative to \file{/api/v1/}, as in \code{GET patients/\{id\}};
the signature gives the full path. \code{\{id\}} is always a whole number, the row of the patient,
encounter, recording, cohort, database or user in the catalog, unless an entry says otherwise.
The parameters of a \code{GET} or \code{DELETE} request are query parameters. Those of
\code{POST}, \code{PUT} and \code{PATCH} requests are the fields of a JSON body, unless the entry
says that the body is a file.

The examples use \cmd{curl} with a cookie file (\file{jar}) against a server on
\code{http://localhost:8090}, the default port, and the synthetic demo data.

\section{Conventions}
\label{sec:http-api-conventions}

\subsection{Signing in and sessions}
\label{sec:http-api-signin}

Every request needs a signed-in user, except \code{GET sign-in} and \code{POST login}. How a
request is signed in depends on the setting \code{sign\_in} (\cref{chap:sign-in}), which
\code{GET sign-in} reports.

\paragraph{Password sign-in} (\code{sign\_in = "password"}, the default). \code{POST login}
checks the user name and the password, and sets the session cookie. The cookie is named
\code{medorama\_<port>} after the server's port (\code{medorama\_8090} by default) and has the
attributes \code{HttpOnly}, \code{SameSite=Strict} and \code{Path=/}, and \code{Secure} when the
setting \code{secure\_cookies} is on (\cref{chap:settings-reference}). It holds a random token of
64 hexadecimal digits; the catalog keeps only the token's hash.

A session ends in any of these cases:
\begin{itemize}
\item after \code{session\_hours} hours (12 by default), the cookie's lifetime;
\item when it has not been used for \code{session\_idle\_minutes} minutes (60 by default; with
  0, a session never ends for this reason). The requests that the page makes by itself,
  \code{GET version} and any request with the header \code{X-Background: 1}, do not count as
  use, so a page left open does not keep its session alive;
\item at \code{POST logout}, when an administrator ends the user's sessions
  (\code{PATCH users/\{id\}} with \code{sign\_out}), or when the user is disabled.
\end{itemize}

After five failed sign-ins with one user name, that name is locked for 60 seconds, and further
attempts are answered \code{429}. Each failed attempt is answered after a pause of one second.
Sign-ins, failed sign-ins and sign-outs everywhere are written to the audit log
(\cref{chap:users-privacy}).

\paragraph{Proxy sign-in} (\code{sign\_in = "proxy"}). An authenticating reverse proxy in front
of \md{} signs the user in and names the user's account in a header (\code{proxy\_user\_header},
by default \code{X-Remote-User}). \md{} believes the header only on connections from
\code{proxy\_addresses}, and keeps no session of its own: every request is signed in by the
header, and the proxy's session is the user's session.
\begin{itemize}
\item A request from an address in \code{proxy\_addresses} whose header names an enabled account
  (compared ignoring case) is signed in as that account.
\item A header that names no account, a disabled account, or two accounts that differ only in
  case is answered \code{403}, with the reason in \code{detail}: \code{your institution's sign-in
  names you NAME, and Medorama has no account of that name: ask an administrator}. The page
  displays this reason.
\item A request without the header, with an empty header, or from any other address is not signed
  in, and is answered \code{401}.
\item \code{POST login} and \code{POST me/password} are refused with \code{400}, and
  \code{POST logout} answers the proxy's sign-out address (\code{proxy\_logout\_url}).
\item The audit log receives one \code{login} entry per visit (detail \code{proxy}), and one
  \code{login\_failed} entry per visit for a name without an account.
\end{itemize}

In both modes, a request without a signed-in user is answered \code{401} with
\code{\{"detail": "sign in required"\}}.

\subsection{The change header}
\label{sec:http-api-header}

Every request under \file{/api/} other than \code{GET} and \code{HEAD}, \code{POST login}
included, must carry the header \code{X-Medorama: 1}. Without it, the answer is \code{403}
\code{\{"detail": "missing X-Medorama header"\}}. The page sends this header. A form or a link on
another site cannot add it, so no other site can make a signed-in browser change anything; this
protects against cross-site request forgery in both sign-in modes. A JSON body is sent with
\code{Content-Type: application/json}.

\subsection{Answers}
\label{sec:http-api-answers}

Answers are JSON (\code{application/json}), except the exports (CSV, JSON or XLSX files sent as
attachments) and the patients' files. Every answer under \file{/api/} carries
\code{Cache-Control: no-store}, except the patients' files, which carry
\code{Cache-Control: private, max-age=300} and \code{Vary: Cookie} (the same address answers
differently in Privacy Mode). Times are Unix seconds (UTC) unless an entry says otherwise.

Every answer also carries the security headers \code{X-Content-Type-Options: nosniff} and
\code{Referrer-Policy: no-referrer}, and a content security policy: the page and the API may not
be framed by another site (\code{frame-ancestors 'none'}, \code{X-Frame-Options: DENY}); the
patients' files, the viewers' pages under \file{/viewers/} and the waveform viewer under
\file{/viewer/} may be framed by \md{}'s own page only (\code{frame-ancestors 'self'},
\code{X-Frame-Options: SAMEORIGIN}). An SVG or XML file is sent with a sandboxing policy, so
that a script it holds cannot run with the rights of the catalog.

\subsection{Errors}
\label{sec:http-api-errors}

An error found by \md{} is a JSON object \code{\{"detail": "..."\}}, the message written for a
person, with one of the status codes of \cref{tab:http-api-status}.

\begin{xltabular}{\linewidth}{@{}P{1.4cm}L@{}}
\caption{The status codes of the API.}
\label{tab:http-api-status}\\
\toprule
{\tabhead Status} & {\tabhead Meaning} \\
\midrule
\endfirsthead
\toprule
{\tabhead Status} & {\tabhead Meaning} \\
\midrule
\endhead
\bottomrule
\endfoot
\code{200} & Done; the answer is the result. \\
\code{202} & An import started in the background (\code{background}): the answer is its task
  (\code{GET import-tasks/\{id\}}). \\
\code{400} & The request is wrong: a missing or invalid value, a rule of the catalog broken (for
  example the last enabled admin disabled), or an action the user's role does not allow where
  the check is made by the catalog rather than by the route (the message says which). \\
\code{401} & No valid session (\code{sign in required}), or a wrong user name or password at
  \code{POST login}. \\
\code{403} & Not allowed: the role (\code{only admins may do this}), the PHI permission, the
  change header missing, or an answer the privacy gate does not send in Privacy Mode. \\
\code{404} & No such patient, encounter, recording, cohort, file, task or route
  (\code{not found}). \\
\code{413} & The body is larger than the route accepts: 1\,GiB for the uploads
  (\code{tables/inspect}, \code{tables}, \code{tables/files}), 2\,MB for any other request. \\
\code{429} & Too many failed sign-ins with this user name: try again in 60 seconds. \\
\code{500} & A fault of the server; the message is also written on its standard error. \\
\code{503} & The catalog is busy saving an import (SQLite has one writer): the answer carries
  \code{Retry-After: 5}, and its \code{detail} says to try again when the import is done.
  \code{POST viewer/start} also answers \code{503} when the viewer cannot be started. \\
\end{xltabular}

A request whose body or parameters cannot be read at all is refused before it reaches \md{}'s
code, with the status of the web framework and a plain-text message instead of a JSON
\code{detail}: \code{400} (a body that is not JSON, a path parameter that is not a number),
\code{415} (a JSON body without \code{Content-Type: application/json}) or \code{422} (a field
missing or of the wrong type). The removals (\code{POST encounters/remove},
\code{POST encounters/remove-matching}, \code{POST recordings/remove},
\code{POST recordings/remove-matching}) also refuse with \code{422} a field they do not know, so
that a misspelled field is never ignored.

\subsection{Roles, the PHI permission and Privacy Mode}
\label{sec:http-api-roles}

Each entry says who may call the route: \emph{any} signed-in user, an \emph{analyst} (the
analyst or admin role), or an \emph{admin}. Some routes also need the PHI permission. Users,
roles and the PHI permission are described in \cref{chap:users-privacy}. Identifiers (names,
dates of birth, MRNs, account numbers) are sent only to users who have the PHI permission and
whose session is out of Privacy Mode.

Privacy Mode is a property of the session. It is always on for users without the PHI
permission. For the others, it is what \code{POST privacy} last chose, kept by the server for
the session (a new sign-in starts in the user's default mode); by default it is on, except for
administrators with the PHI permission, who start out of it. In Privacy Mode, every answer passes
the privacy gate (\cref{chap:protecting}):
\begin{itemize}
\item every JSON answer is rewritten on its way out: identifiers become placeholders
  (\code{[Patient Name]}, \ldots) and pseudonyms (\code{MO-...}), dates move back by the
  offset of the settings (on a patient's own routes, \code{patients/\{id\}/...}, by that
  patient's own offset), ages above 89 become 90+. The answers about the session itself
  (\code{login}, \code{logout}, \code{me}, \code{privacy}, \code{sign-in}) are left as they
  are;
\item the text of notes and reports is sent by default, with the patient's own names and date of
  birth replaced, and marked with a warning (it may still name the patient or others); it is
  \code{[Withheld]} only when the admin withholds \emph{The text of notes and reports}
  (\code{notes} in \code{PUT admin/deid});
\item a file is sent only as the gate allows (\cref{sec:http-api-files});
\item any other answer that carries data and that the gate has not cleared is refused, with
  \code{403} \code{\{"detail": "this answer is not available in Privacy Mode"\}}, so that a
  route that forgets the gate cannot leak.
\end{itemize}
The exports (\code{databases/\{id\}/export}, \code{cohorts/\{id\}/export}) pass the gate in
Privacy Mode and also when asked with \code{privacy=1}.

\subsection{The patients' files}
\label{sec:http-api-files}

Out of Privacy Mode, a file is sent as it is. In Privacy Mode, the privacy gate decides by the
file's kind (\cref{sec:protecting-files}):
\begin{itemize}
\item a WFDB header (\file{.hea}): de-identified, its record and file names replaced by tokens;
\item a WFDB signal file (\file{.dat}): as it is (it holds only samples);
\item a WFDB signal file in a MATLAB file (\file{.mat}): its first matrix only, its descriptive
  text and the matrix's name blanked; one that is not a MAT file WFDB reads is refused
  (\code{403});
\item a GE MUSE XML ECG (\file{.xml} with \code{<RestingECG>}): de-identified; any other XML
  file is refused (\code{403});
\item a DICOM file: its header rewritten (the pseudonym, the dates moved, the UIDs replaced),
  its pixels as they are;
\item an EDF, EDF+, BDF or BDF+ file (\file{.edf}, \file{.bdf}): its header de-identified
  (the patient field \code{X X X X}; in EDF+ and BDF+, the recording field
  \code{Startdate <moved date> X X X}, in plain EDF and BDF blank; the start date and time moved
  back), the texts of EDF+ annotations removed, the samples as they are. A file with
  annotation signals is rewritten in memory, so such a file larger than 64\,MB is refused
  (\code{403}; \cref{sec:protecting-edf});
\item a picture, a video or a PDF (\file{.png}, \file{.jpg}, \file{.jpeg}, \file{.gif},
  \file{.webp}, \file{.bmp}, \file{.svg}, \file{.mp4}, \file{.m4v}, \file{.webm}, \file{.mov},
  \file{.ogv}, \file{.pdf}): as it is;
\item anything else, JSON ECGs, text files and signal files with an extension other than
  \file{.dat} or \file{.mat} among them: refused (\code{403}).
\end{itemize}
A file whose content the gate cannot inspect (a picture, a video, a PDF, DICOM pixels) carries
the header \code{x-medorama-may-contain-phi: 1}, on which the page shows its warning. In Privacy
Mode the files of an attachment are named by tokens, \code{mo-f} followed by the file's number
in four digits and its extension (\code{mo-f0003.dcm}), never by their names.

\subsection{Pages and the viewers}
\label{sec:http-api-pages}

Besides the API, the server answers:
\begin{itemize}
\item \code{GET /}: the page; \code{GET /static/...}: its scripts and styles;
\item \code{GET /viewers/ecg/...} and \code{GET /viewers/pacs/...}: the pages of the ECG and
  PACS viewers (\cref{chap:viewer-api}), which hold no data: they read the patients' files
  through \code{patients/\{id\}/file/...} with the user's session;
\item \file{/viewer/}: the VitalStream viewer, behind the catalog's sign-in and its gateway,
  which passes only the requests of the viewer API that the page needs, each with its rule
  (\cref{chap:viewer-api}). \code{/viewer} redirects to \code{/viewer/}; a page request without
  a session redirects to the catalog's page. The gateway answers \code{413} to a request body
  over 1\,MiB, \code{415} to a \code{POST} not sent as JSON, and \code{403} to a \code{POST}
  without \code{X-Medorama: 1}.
\end{itemize}
Any other path under \file{/api/v1/} is answered \code{404}.

\section{The session and the settings}
\label{sec:http-api-session}

\apientry{GET sign-in}{How users sign in}
\begin{lstlisting}[style=signature]
GET /api/v1/sign-in
\end{lstlisting}
Tells the page which sign-in method the catalog uses (\cref{chap:sign-in}), so that it shows the
password form or explains that the institution's sign-in names the user. It needs no session and
passes the privacy layer unchanged.
\returns{\code{method}: \code{password} or \code{proxy}; \code{logout\_url}: the proxy's sign-out
address (\code{proxy\_logout\_url}), empty in password mode or when none is set.}
\begin{lstlisting}[style=shell]
curl -s http://localhost:8090/api/v1/sign-in
\end{lstlisting}
\begin{lstlisting}[language=JSON]
{"method": "proxy", "logout_url": "/oauth2/sign_out"}
\end{lstlisting}

\apientry{POST login}{Sign in}
\begin{lstlisting}[style=signature]
POST /api/v1/login   {"username": NAME, "password": PASSWORD}
\end{lstlisting}
Checks the user name and the password of an enabled user, starts a session and sets its cookie
(\cref{sec:http-api-signin}). It needs no session, but it needs the change header. Password
sign-in only. Audited (\code{login}, or \code{login\_failed} with the name tried).
\begin{params}
\item[username] The user name.
\item[password] The password.
\end{params}
\returns{\code{\{"user": USER\}}: the user (\code{id}, \code{username}, \code{display\_name},
\code{role}, \code{phi}, \code{disabled}, \code{phi\_allowed}, \code{privacy},
\code{privacy\_notes}), and the header \code{Set-Cookie}.}
\raises{\code{400} proxy sign-in (\code{this catalog signs users in through the institution's
sign-in, not with a password}); \code{401} a wrong user name or password; \code{403} the change
header missing; \code{429} too many failed attempts.}
\begin{lstlisting}[style=shell]
curl -s -c jar -H 'X-Medorama: 1' -H 'Content-Type: application/json' \
  -d '{"username": "admin", "password": "correct-horse-battery"}' \
  http://localhost:8090/api/v1/login
\end{lstlisting}

\apientry{POST logout}{Sign out}
\begin{lstlisting}[style=signature]
POST /api/v1/logout   [{"everywhere": BOOL}]
\end{lstlisting}
Ends this session and removes the session cookie and the Privacy Mode cookie, so that the next
sign-in starts in the user's default mode. The body is optional. The route needs no valid
session: without one, it only removes the cookies. With proxy sign-in, \md{} keeps no session:
the route only forgets the Privacy Mode choice, and the answer names the proxy's sign-out, to
which the page then goes.
\begin{params}
\item[everywhere] \code{true}: end every session of the user, on every device (password sign-in
  only). Audited (\code{signed\_out\_everywhere}, with the number of sessions). Default
  \code{false}.
\end{params}
\returns{\code{\{"logout\_url": URL\}}: empty with password sign-in; with proxy sign-in, the
setting \code{proxy\_logout\_url}.}

\apientry{GET me}{Who is signed in}
\begin{lstlisting}[style=signature]
GET /api/v1/me
\end{lstlisting}
The session's user, in the session's mode: in Privacy Mode \code{phi} is \code{false} and
\code{privacy} \code{true}; \code{phi\_allowed} says whether the user may leave Privacy Mode.
Any user.
\returns{\code{\{"user": USER\}}, as for \code{POST login}.}
\begin{lstlisting}[language=JSON]
{"user": {"id": 2, "username": "reader1", "display_name": "Reader One", "role": "viewer",
          "phi": false, "disabled": false, "phi_allowed": false, "privacy": true,
          "privacy_notes": false}}
\end{lstlisting}

\apientry{POST me/password}{Change one's password}
\begin{lstlisting}[style=signature]
POST /api/v1/me/password   {"old": PASSWORD, "new": PASSWORD}
\end{lstlisting}
Changes the signed-in user's password. Any user, with password sign-in only. Audited
(\code{password\_changed}).
\begin{params}
\item[old] The current password.
\item[new] The new password: at least 10 characters.
\end{params}
\returns{\code{\{\}}.}
\raises{\code{400} the new password is too short, or proxy sign-in (\code{this catalog signs
users in through the institution's sign-in: Medorama's passwords are not used}); \code{403} the
current password is wrong.}

\apientry{POST privacy}{Switch Privacy Mode}
\begin{lstlisting}[style=signature]
POST /api/v1/privacy   {"on": BOOL}
\end{lstlisting}
Switches the session's Privacy Mode. The choice is kept by the server for the session, never
in the browser: in password mode until the session ends, in proxy mode for the visit (at most
\code{session\_hours}); a new sign-in starts in the user's default mode. Any user may switch it
on; switching it off needs the PHI permission. Audited (\code{privacy\_on},
\code{privacy\_off}).
\begin{params}
\item[on] \code{true}: Privacy Mode on; \code{false}: off.
\end{params}
\returns{\code{\{"privacy": BOOL\}}.}
\raises{\code{403} switching it off without the PHI permission.}

\apientry{GET config}{The server's settings for the page}
\begin{lstlisting}[style=signature]
GET /api/v1/config
\end{lstlisting}
What the page needs to know about the server. Any user.
\returns{\code{version} (of \md{}), \code{timezone}, \code{viewer\_url} (a viewer running
elsewhere, opened in a new browser tab; empty when \md{} runs its own), \code{auto\_index}
(whether new and changed recordings are indexed while the server runs), \code{workers} (the
number of job workers), \code{user} (as \code{GET me}), \code{sign\_in} (\code{password} or
\code{proxy}; the page hides \ui{Change password}, \ui{Sign out everywhere} and the password
fields of \ui{Admin} $\rightarrow$ \ui{Users} with \code{proxy}), \code{results\_dir} and \code{index\_dir} (the
folders, for admins only; \code{null} for the others), \code{viewer} (as \code{GET viewer}),
\code{deid} (as \code{GET deid}; for users who are not admins, without \code{offset},
\code{offset\_seconds} and \code{spread\_days}) and \code{display} (the date and time formats, as
\code{PUT admin/display} answers).}

\apientry{GET deid}{What Privacy Mode withholds}
\begin{lstlisting}[style=signature]
GET /api/v1/deid
\end{lstlisting}
The privacy gate's settings: what Privacy Mode withholds and how it moves dates. Admins: with
the offset, any date shown would give the real one, so other users learn what is withheld from
\code{GET config}, without the offset.
\returns{\code{fields} (the keys of the fields withheld), \code{offset}
(\code{\{years, months, days, hours\}}), \code{offset\_seconds}, \code{spread\_days}, and
\code{choices}: every field that can be withheld, each with its \code{key}, \code{label},
\code{hipaa} and \code{default}. The keys are \code{names}, \code{birth\_date}, \code{dates},
\code{age\_over\_89}, \code{mrn}, \code{account\_numbers} and \code{notes}; all but
\code{notes} are withheld by default.}
\begin{lstlisting}[language=JSON]
{"fields": ["account_numbers", "age_over_89", "birth_date", "dates", "mrn", "names"],
 "offset": {"years": 0, "months": 0, "days": 36570, "hours": 0},
 "offset_seconds": 3159648000, "spread_days": 365, "choices": [...]}
\end{lstlisting}

\apientry{PUT admin/deid}{Change what Privacy Mode withholds}
\begin{lstlisting}[style=signature]
PUT /api/v1/admin/deid   {"fields": [KEY...], "offset": {...}, "spread_days": N}
\end{lstlisting}
Saves the privacy gate's settings (\cref{chap:users-privacy}). Admins. Audited
(\code{deid\_settings}). Unknown field keys are dropped; a field left out of the body takes
its default.
\begin{params}
\item[fields] The keys of the fields to withhold (as \code{GET deid} lists them).
\item[offset] How far every date moves back: \code{years}, \code{months}, \code{days},
  \code{hours}, each 0 or more (at most 1000 years, 12\,000 months, 400\,000 days,
  10\,000\,000 hours). Default 36\,570 days.
\item[spread\_days] Each patient's own extra days, drawn once per patient, from 0 to this
  (30 to 3650; default 365).
\end{params}
\returns{the new settings, as \code{GET deid}.}
\raises{\code{400} a value out of range; \code{403} a user who is not an admin.}

\apientry{PUT admin/display}{How dates and times are written}
\begin{lstlisting}[style=signature]
PUT /api/v1/admin/display   {"date": FORMAT, "clock": CLOCK}
\end{lstlisting}
How the page writes dates and times, for every user. Admins. Audited (\code{display\_changed}).
\begin{params}
\item[date] \code{ymd} (2024-01-31), \code{mdy} (01/31/2024), \code{dmy} (31/01/2024) or
  \code{dmy\_dot} (31.01.2024).
\item[clock] \code{24h} (14:05) or \code{12h} (2:05 PM).
\end{params}
\returns{\code{\{date, clock, dates, clocks\}}: the setting and the choices (each with its
\code{key}, \code{example} and \code{label}).}
\raises{\code{400} an unknown format; \code{403} a user who is not an admin.}

\apientry{GET admin/addons}{The installed add-ons}
\begin{lstlisting}[style=signature]
GET /api/v1/admin/addons
\end{lstlisting}
The adapters installed in the adapter folder (\cref{chap:addons}). Installing is done on the
server only, with \cmd{medorama addons install} (\cref{chap:cli}). Admins.
\returns{\code{\{"dir": FOLDER, "items": [\{name, version, implementation, program,
formats\}]\}}.}

\apientry{GET version}{The data's version}
\begin{lstlisting}[style=signature]
GET /api/v1/version
\end{lstlisting}
A number that moves at every import, removal or change of databases, patients, encounters or
recordings: open pages ask for it and redraw when it moves. Does not count as use of the
session. Any user.
\returns{\code{\{"data": N\}}.}

\section{Patients}
\label{sec:http-api-patients}

\apientry{GET patients}{Find patients}
\begin{lstlisting}[style=signature]
GET /api/v1/patients?q=&has=&cohort=&mark=&limit=&offset=
\end{lstlisting}
The patients, latest encounter first. Any user; identifiers only as \cref{sec:http-api-roles}
says.
\begin{params}
\item[q] The search text (\cref{chap:patients}).
\item[has] Categories the patients must have, comma-separated (\code{icu,ecg}); \code{!icu}:
  without that category.
\item[cohort] Only the members of this cohort (one the user may see).
\item[mark] Each item says whether the patient is in this cohort (\code{member}).
\item[limit, offset] The page of the list: \code{limit} 1 to 1000 (default 100),
  \code{offset} default 0.
\end{params}
\returns{\code{\{total, items, categories\}}: each item has \code{id}, \code{patient\_id},
\code{pseudonym}, \code{sex}, \code{age}, \code{encounters}, \code{latest}, \code{latest\_where},
\code{kinds}, \code{items}, \code{sources}, and \code{phi} (the identifiers) for users who see
them; \code{categories} lists the categories present, each as \code{[key, label]}.}
\raises{\code{404} a cohort the user may not see.}
\begin{lstlisting}[style=shell]
curl -s -b jar 'http://localhost:8090/api/v1/patients?has=icu&limit=20'
\end{lstlisting}

\apientry{GET patients/\{id\}}{A patient's summary}
\begin{lstlisting}[style=signature]
GET /api/v1/patients/{id}
\end{lstlisting}
A patient's summary: the identity (protected health information), the encounters with their
bed stays, problems, medications and allergies; the stays found in the connected databases
(\code{visits}) and the connected databases that did not answer (\code{unavailable}). Any user;
audited (\code{patient\_view}) for a user who sees identifiers.
In Privacy Mode, the patient's own names and date of birth are also replaced wherever a text
of the answer names them.

\apientry{GET patients/\{id\}/journey}{A patient's journey}
\begin{lstlisting}[style=signature]
GET /api/v1/patients/{id}/journey
\end{lstlisting}
The journey of \cref{chap:patients}: its columns, its cells of six hours (their items and the
worst level among them), and its spans (stays, drugs, recordings). Any user.

\apientry{GET patients/\{id\}/report}{A patient's report, or a search in it}
\begin{lstlisting}[style=signature]
GET /api/v1/patients/{id}/report?q=
\end{lstlisting}
The entries of the patient's report. With \code{q}, the search in it: counts per category,
the words found highlighted, the fuzzy matches. Any user. It reads the patient's own items
only, and in Privacy Mode their text as the session is shown it. A search by a user who sees
identifiers is audited (\code{phi\_search}), and one in Privacy Mode too
(\code{privacy\_search}).
\begin{params}
\item[q] The words to find (optional).
\end{params}

\apientry{GET patients/\{id\}/detail}{One entry's detail}
\begin{lstlisting}[style=signature]
GET /api/v1/patients/{id}/detail?key=KEY
\end{lstlisting}
The detail of one entry of the report: the rows of a panel, a text, an attachment, the ICU
overview, the notes around it. Any user.
\begin{params}
\item[key] The entry's key, as the report gave it (in Privacy Mode, with its time moved back:
  the server moves it forward again).
\end{params}

\apientry{GET patients/\{id\}/trends}{A patient's numeric series}
\begin{lstlisting}[style=signature]
GET /api/v1/patients/{id}/trends
\end{lstlisting}
Every numeric series of the patient (results, vital signs), for the trends. Any user.

\apientry{GET patients/\{id\}/file/\{item\}}{An item's attached file}
\begin{lstlisting}[style=signature]
GET /api/v1/patients/{id}/file/{item}
GET /api/v1/patients/{id}/file/{item}/
\end{lstlisting}
The file attached to an item of the patient (\code{item} is the item's number), such as a
picture, a PDF document, a text file, an ECG (GE MUSE XML, DICOM, EDF) or a WFDB header. For an
attached folder, such as a DICOM series, the answer lists its files instead. In Privacy Mode, the
file passes the gate (\cref{sec:http-api-files}). Any user.
\returns{the file, or for a folder \code{\{"files": [NAME...]\}}; in Privacy Mode the folder
lists only the files the gate would send, by their tokens (\code{mo-f0003.dcm}).}
\raises{\code{404} no such item or file; \code{403} a file the gate does not send in Privacy
Mode.}

\apientry{GET patients/\{id\}/file/\{item\}/\{name\}}{A file of an attachment}
\begin{lstlisting}[style=signature]
GET /api/v1/patients/{id}/file/{item}/{name}
\end{lstlisting}
A file of an item's attachment by its name: the attached file itself, the signal files that
its WFDB header names, or a file of an attached folder; nothing else of the folder in which the
attachment lies. In Privacy Mode, by its token only. Any user.
\raises{\code{404} a name that is not one of the attachment's files (or not its token in
Privacy Mode); \code{403} as above.}
\begin{lstlisting}[style=shell]
curl -s -b jar -o mo-f0000.hea 'http://localhost:8090/api/v1/patients/3/file/41/mo-f0000.hea'
\end{lstlisting}

\section{Encounters and recordings}
\label{sec:http-api-encounters}

\apientry{GET encounters}{Find encounters}
\begin{lstlisting}[style=signature]
GET /api/v1/encounters?q=&unit=&from=&to=&status=&limit=&offset=
\end{lstlisting}
The encounters that match a search. Any user.
\begin{params}
\item[q] The search text: words and \code{field:value} terms (\cref{chap:patients}).
\item[unit] Only stays in this unit.
\item[from, to] Only stays that overlap \code{[from, to]} (Unix seconds).
\item[status] The summary status of the encounter's recordings (\code{ready}, \code{new},
  \ldots, or \code{none}: without recordings).
\item[limit, offset] The page of the list.
\end{params}
\returns{\code{\{total, items\}}: each item with \code{id}, \code{patient},
\code{encounter\_id}, \code{patient\_id}, \code{sex}, \code{admit\_time},
\code{discharge\_time}, \code{admit\_unit}, \code{disposition}, \code{units},
\code{recordings}, \code{status}, \code{status\_counts}, \code{sources}, and the identifiers for
users who see them.}

\apientry{GET encounters/\{id\}}{An encounter}
\begin{lstlisting}[style=signature]
GET /api/v1/encounters/{id}
\end{lstlisting}
An encounter with its bed stays, its recordings (their status, their link, their results: the
newest job of each action other than \code{index} and \code{identify}, with its summary), the counts of its items and the patient's other
encounters. Adds it to the user's history (\code{GET history}). Any user; audited
(\code{phi\_view}) for a user who sees identifiers.
\raises{\code{404} no such encounter.}

\apientry{GET encounters/\{id\}/ehr}{An encounter's items}
\begin{lstlisting}[style=signature]
GET /api/v1/encounters/{id}/ehr?domain=&patient_level=&limit=&offset=
\end{lstlisting}
The encounter's items in time order. Any user.
\begin{params}
\item[domain] Only the items of this domain (category).
\item[patient\_level] \code{true}: the patient's items that belong to no encounter instead.
\item[limit, offset] The page: \code{limit} 1 to 5000 (default 500), \code{offset} default 0.
\end{params}
\returns{\code{\{total, items\}}.}

\apientry{POST encounters/remove}{Remove encounters}
\begin{lstlisting}[style=signature]
POST /api/v1/encounters/remove   {"ids": [ID...], "dry_run": BOOL}
\end{lstlisting}
Removes encounters with their bed stays and their items; the recordings stay (unlinked).
Encounters hold no files or results, so there is nothing to purge. Admins. Audited
(\code{remove\_encounter}, each encounter) unless a dry run.
\begin{params}
\item[ids] The encounters' numbers.
\item[dry\_run] \code{true}: say what would be removed, and change nothing.
\end{params}
\returns{\code{\{removed, report\}}: the number of encounters removed and the report
(\code{\{dry\_run, encounters, patients, bed\_stays, ehr\_items, ...\}}).}
\raises{\code{422} any other field (\code{purge} among them).}

\apientry{POST encounters/remove-matching}{Remove the encounters a search finds}
\begin{lstlisting}[style=signature]
POST /api/v1/encounters/remove-matching   {"search": {...}, "dry_run": BOOL}
\end{lstlisting}
Removes the encounters that a search finds, as \code{POST encounters/remove}. Admins. Audited
(\code{remove\_matching}) unless a dry run.
\begin{params}
\item[search] The search, with the fields of \code{GET encounters} (\code{q}, \code{unit},
  \code{from}, \code{to}, \code{status}).
\item[dry\_run] \code{true}: say what would be removed, and change nothing.
\end{params}
\returns{the report, as in \code{POST encounters/remove}: \code{\{dry\_run, encounters,
patients, bed\_stays, ehr\_items, ...\}}.}
\raises{\code{422} any other field: the body takes \code{search} and \code{dry\_run} only, so that
a field the route does not read (\code{q}, \code{unlinked}) cannot widen the removal to every
encounter.}

\apientry{GET recordings}{Find recordings}
\begin{lstlisting}[style=signature]
GET /api/v1/recordings?q=&unlinked=&limit=&offset=
\end{lstlisting}
The recordings in the catalog. Any user.
\begin{params}
\item[q] The search text.
\item[unlinked] \code{true}: only the recordings linked to no encounter.
\item[limit, offset] The page: \code{limit} 1 to 1000 (default 100), \code{offset} default 0.
\end{params}
\returns{\code{\{total, items\}}.}

\apientry{GET recordings/\{id\}}{A recording}
\begin{lstlisting}[style=signature]
GET /api/v1/recordings/{id}
\end{lstlisting}
A recording, the encounters it is linked to and the results of the actions run on it. Any user.
\returns{\code{\{recording, encounters, results\}}: each encounter with \code{id},
\code{encounter\_id} and \code{method} (how it was linked).}
\raises{\code{404} no such recording.}

\apientry{POST recordings/remove}{Remove recordings}
\begin{lstlisting}[style=signature]
POST /api/v1/recordings/remove   {"ids": [ID...], "purge": BOOL, "dry_run": BOOL}
\end{lstlisting}
Removes recordings from the catalog. The recording files are never touched. Admins. Audited
(\code{remove\_recording}, each recording) unless a dry run.
\begin{params}
\item[ids] The recordings' numbers.
\item[purge] \code{true}: also delete their results and indexes (inside \code{results\_dir} and
  \code{index\_dir}).
\item[dry\_run] \code{true}: say how many would go, and change nothing.
\end{params}
\returns{\code{\{"dry\_run": BOOL, "removed": N\}}.}
\raises{\code{422} any other field.}

\apientry{POST recordings/remove-matching}{Remove the recordings a search finds}
\begin{lstlisting}[style=signature]
POST /api/v1/recordings/remove-matching   {"q": TEXT, "unlinked": BOOL, "purge": BOOL,
                                           "dry_run": BOOL}
\end{lstlisting}
Removes the recordings that a search finds, as \code{POST recordings/remove}. Admins. Audited
(\code{remove\_matching}) unless a dry run.
\begin{params}
\item[q, unlinked] The search, as in \code{GET recordings}.
\item[purge] As in \code{POST recordings/remove}.
\item[dry\_run] \code{true}: say how many would go, and change nothing.
\end{params}
\returns{\code{\{"dry\_run": BOOL, "recordings": N\}}.}
\raises{\code{422} any other field.}

\apientry{POST recordings/\{id\}/link}{Link a recording manually}
\begin{lstlisting}[style=signature]
POST /api/v1/recordings/{id}/link   {"encounter": ID | null}
\end{lstlisting}
Links a recording to an encounter by hand (\cref{chap:administration}). Admins.
\begin{params}
\item[encounter] The encounter's number; \code{null}: forget the recording's manual links.
\end{params}
\returns{\code{\{\}}.}

\apientry{POST recordings/\{id\}/decision}{Decide on a proposed link}
\begin{lstlisting}[style=signature]
POST /api/v1/recordings/{id}/decision   {"encounter_id": TEXT, "decision": "link" | "never" | null}
\end{lstlisting}
Decides on the link of a recording to an encounter that the catalog proposes. Admins. Audited
(\code{link\_decision}).
\begin{params}
\item[encounter\_id] The encounter's identifier (as imported).
\item[decision] \code{link}: link them; \code{never}: never link them; \code{null}: forget the
  decision.
\end{params}
\returns{\code{\{\}}.}

\apientry{POST check}{Check the recordings again}
\begin{lstlisting}[style=signature]
POST /api/v1/check
\end{lstlisting}
Checks every recording again: whether its files are still there, whether they changed, and
whether its index folder is still there. Admins. Audited
(\code{check}).
\returns{the counts of the check.}

\section{Review}
\label{sec:http-api-review}

\apientry{GET review}{What waits for a decision}
\begin{lstlisting}[style=signature]
GET /api/v1/review
\end{lstlisting}
The Review tab (\cref{chap:administration}): the rows held back at import, the items waiting for
their patient or encounter; for users who see identifiers also the possible duplicate patients
and the recordings whose files disagree with the catalog. Admins.

\apientry{POST review/held/\{id\}}{Decide on a held-back row}
\begin{lstlisting}[style=signature]
POST /api/v1/review/held/{id}   {"action": "apply" | "dismiss"}
\end{lstlisting}
Decides on a row held back at import. Admins. Audited (\code{apply\_held\_row},
\code{dismiss\_held\_row}).
\begin{params}
\item[action] \code{apply}: apply the row, with the fix its problem calls for;
  \code{dismiss}: drop it.
\end{params}
\returns{\code{\{\}}.}
\raises{\code{400} another action.}

\apientry{POST review/merge}{Merge two patients}
\begin{lstlisting}[style=signature]
POST /api/v1/review/merge   {"a": PATIENT_ID, "b": PATIENT_ID}
\end{lstlisting}
Merges patient \code{a} into patient \code{b}: \code{a}'s data become \code{b}'s, and
the recordings and waiting items are linked again. Admins with the PHI permission, out of
Privacy Mode. Audited (\code{merge\_patients}).
\begin{params}
\item[a, b] The patients' identifiers (\code{patient\_id}, as imported).
\end{params}
\returns{\code{\{\}}.}
\raises{\code{403} a user who is not an admin, without the PHI permission, or in Privacy Mode;
\code{400} \code{a} and \code{b} the same patient.}

\apientry{POST review/distinct}{Two patients are different people}
\begin{lstlisting}[style=signature]
POST /api/v1/review/distinct   {"a": PATIENT_ID, "b": PATIENT_ID}
\end{lstlisting}
Records that two patients proposed as duplicates are different people, so that they are no
longer proposed. Admins with the PHI permission, out of Privacy Mode. Audited
(\code{distinct\_patients}).
\returns{\code{\{\}}.}
\raises{\code{403} a user who is not an admin, without the PHI permission, or in Privacy Mode.}

\section{Databases and imports}
\label{sec:http-api-databases}

\apientry{GET databases}{The databases in use}
\begin{lstlisting}[style=signature]
GET /api/v1/databases
\end{lstlisting}
The databases of the Databases tab (\cref{chap:databases-tab}): their rows, the patients they
reach, their categories, their held-back rows and waiting items; for a connected database its
\code{connection}. Any user.
\returns{\code{\{"items": [...]\}}.}

\apientry{GET databases/\{id\}/rows}{A database's rows}
\begin{lstlisting}[style=signature]
GET /api/v1/databases/{id}/rows?q=&f<pos>=&sort=&desc=&limit=&offset=
\end{lstlisting}
One page of a database's rows, as imported: its columns (each with its \code{header} and the
catalog field it \code{maps\_to}), the columns not shown, and each row's cells, status, patient,
entry and link to its file. Any user; in Privacy Mode through the gate, and the search, the
filters and the sort then use the cells as the session is shown them (a sort by a placeholder
column is ignored); such a search or filter is audited (\code{privacy\_search}). A sort puts
numbers in numeric order before text, and empty cells last. Audited (\code{view\_database}) for
a user who sees identifiers.
\begin{params}
\item[q] Text to find in the rows.
\item[f<pos>] A filter on the column at position \code{<pos>} (\code{f0}, \code{f3}, \ldots).
\item[sort] The position of the column to sort by; \code{desc=1} (or \code{true}): descending.
\item[limit, offset] The page (default 100 rows from 0).
\end{params}

\apientry{GET databases/\{id\}/export}{A database as CSV}
\begin{lstlisting}[style=signature]
GET /api/v1/databases/{id}/export?privacy=
\end{lstlisting}
The database as a CSV file, as imported. Admins. Through the gate in Privacy Mode, or with
\code{privacy=1} (\code{deid=1} is the earlier name of the same parameter). Audited
(\code{phi\_export} for a user with the PHI permission out of Privacy Mode, else
\code{export\_database}).
\returns{a \code{text/csv} attachment.}

\apientry{GET formats}{The import formats}
\begin{lstlisting}[style=signature]
GET /api/v1/formats
\end{lstlisting}
The import formats the catalog knows: those built in and those described in the site's folders
(\code{formats\_dirs}), the files there that could not be read, and the categories
(\cref{chap:importing}). Admins.
\returns{\code{\{items, problems, categories\}}: each category with its \code{key} and
\code{label}.}

\apientry{POST tables/inspect}{Look at a table before importing it}
\begin{lstlisting}[style=signature]
POST /api/v1/tables/inspect?name=FILE&format=   (body: the file)
\end{lstlisting}
Reads an uploaded table without importing it: its headers, the mapping and the category they
suggest, the format it is in (converted), the catalog's columns by group; for users who see
identifiers, also its first rows. Admins. The body is the file (at most 1\,GiB).
\begin{params}
\item[name] The file's name.
\item[format] \code{auto} (the default: recognized from the file), \code{none} (read as it is),
  or a format's key.
\end{params}
\begin{lstlisting}[style=shell]
curl -s -b jar -H 'X-Medorama: 1' --data-binary @labs.csv \
  'http://localhost:8090/api/v1/tables/inspect?name=labs.csv'
\end{lstlisting}

\apientry{POST tables}{Import a table}
\begin{lstlisting}[style=signature]
POST /api/v1/tables?name=&batch=&source=&category=&mapping=&format=&dry_run=  (body: the file)
\end{lstlisting}
Imports one file as a database, in one batch (\cref{chap:importing}). Admins. The body is the
file (at most 1\,GiB). Audited (\code{import\_table}) unless a dry run.
\begin{params}
\item[name] The file's name (for the import log, and the default batch name).
\item[batch] The database's (batch's) name; default the file's name.
\item[source] The source label.
\item[category] The category of the items when the table has no category column.
\item[mapping] The column mapping, as JSON or TOML text: each catalog column and the header
  that holds it; in TOML at the top level or under \code{[columns]}. Empty: the columns are
  recognized by their headers. With a mapping, a header is read as a catalog column only if the
  mapping names it, or if it is that column's own name and the mapping names no other header
  for it (\cref{chap:importing}).
\item[format] \code{auto} (default), \code{none}, or a format's key.
\item[dry\_run] \code{true}: read and check everything, keep nothing.
\end{params}
\returns{the import's report.}

\apientry{POST tables/files}{Import several files as one database}
\begin{lstlisting}[style=signature]
POST /api/v1/tables/files?batch=&source=&category=&mapping=&format=&dry_run=&background=
     (multipart body: one part named "file" per file)
\end{lstlisting}
Imports one database from one or several files with the same columns, sent as a multipart form
(each part named \code{file}, held in memory, never written to disk). The parameters are those
of \code{POST tables}; \code{batch}, the database's name, is required. Admins.
\begin{params}
\item[background] \code{true}: run in the background; the answer, \code{202}, is the import's
  task (\code{GET import-tasks/\{id\}}).
\end{params}
\returns{the import's report, or with \code{background} its task.}
\raises{\code{400} no \code{batch}, or no file.}
\begin{lstlisting}[style=shell]
curl -s -b jar -H 'X-Medorama: 1' -F file=@vitals-1.csv -F file=@vitals-2.csv \
  'http://localhost:8090/api/v1/tables/files?batch=vitals&background=true'
\end{lstlisting}

\apientry{POST tables/clif}{Import a CLIF folder}
\begin{lstlisting}[style=signature]
POST /api/v1/tables/clif   {"dir": FOLDER, "source": TEXT, "dry_run": BOOL, "background": BOOL}
\end{lstlisting}
Imports a CLIF folder on the server, each table as a database (\cref{chap:clif}). The folder
must lie inside \code{data\_roots}. Admins.
\begin{params}
\item[dir] The folder, on the server.
\item[source] The source label; default \code{clif}.
\item[dry\_run] \code{true}: keep nothing.
\item[background] As in \code{POST tables/files}.
\end{params}
\returns{\code{\{tables, skipped\}}: each table imported with its \code{name}, \code{files} and
\code{report}; the files skipped.}
\raises{\code{400} no such folder, or a folder outside \code{data\_roots}.}

\apientry{GET import-tasks}{The imports in the background}
\begin{lstlisting}[style=signature]
GET /api/v1/import-tasks
\end{lstlisting}
The imports running in the background, and those that ended within the last hour. Admins.
\returns{\code{\{"items": [TASK...]\}}, each task as \code{GET import-tasks/\{id\}}.}

\apientry{GET import-tasks/\{id\}}{An import's progress}
\begin{lstlisting}[style=signature]
GET /api/v1/import-tasks/{id}
\end{lstlisting}
One import's progress, or its end. \code{id} is the task's identifier (24 hexadecimal
digits). Admins.
\returns{\code{\{id, label, state, started, progress, report, error\}}: \code{state} is
\code{running}, \code{done}, \code{cancelled} or \code{failed}; \code{progress} is
\code{\{phase, rows, done, elapsed\_s, rate, left\_s, cancelling\}}, with \code{phase}
\code{rows}, \code{index}, \code{match} or \code{done}, \code{done} the fraction done
(0 to 1), \code{rate} rows per second and \code{left\_s} the seconds left (when it can be
estimated); \code{report} the import's report once done, \code{error} the message once failed or
cancelled.}
\raises{\code{404} no such import.}
\begin{lstlisting}[language=JSON]
{"id": "5f0c2a9e41d7b3c8a6e2f910", "label": "vitals", "state": "running", "started": 1767225600.0,
 "progress": {"phase": "rows", "rows": 120000, "done": 0.42, "elapsed_s": 8.0,
              "rate": 15000.0, "left_s": 11.0, "cancelling": false},
 "report": null, "error": null}
\end{lstlisting}

\apientry{POST import-tasks/\{id\}/cancel}{Stop an import}
\begin{lstlisting}[style=signature]
POST /api/v1/import-tasks/{id}/cancel
\end{lstlisting}
Stops an import in the background, at its next chunk of rows; nothing of it is kept. Audited
(\code{cancel\_import}). Admins.
\returns{the task, as \code{GET import-tasks/\{id\}}.}
\raises{\code{404} no such import.}

\apientry{GET imports}{The batches}
\begin{lstlisting}[style=signature]
GET /api/v1/imports
\end{lstlisting}
The import batches. Admins.
\returns{\code{\{"items": [...]\}}.}

\apientry{POST imports/\{id\}/remove}{Remove a batch}
\begin{lstlisting}[style=signature]
POST /api/v1/imports/{id}/remove   {"purge": BOOL, "dry_run": BOOL}
\end{lstlisting}
Removes an import batch and what it added: the patients, encounters and recordings that
nothing else holds go, and the rows of the database as imported. Admins. Audited
(\code{remove\_import}) unless a dry run.
\begin{params}
\item[purge] \code{true}: also delete the results folders of the actions run on the
  recordings that go with it.
\item[dry\_run] \code{true}: only count.
\end{params}
\returns{the counts of what was (or would be) removed.}

\apientry{GET sources}{The sources}
\begin{lstlisting}[style=signature]
GET /api/v1/sources
\end{lstlisting}
The sources (the labels the imports gave) with what each contributes. Any user.
\returns{\code{\{items, waiting\}}: \code{items} the sources with what they contribute,
\code{waiting} the number of items waiting for their encounter.}

\apientry{POST sources/remove}{Remove a source}
\begin{lstlisting}[style=signature]
POST /api/v1/sources/remove   {"source": LABEL, "purge": BOOL, "dry_run": BOOL}
\end{lstlisting}
Removes everything a source added. Admins.
\begin{params}
\item[source] The source's label (required).
\item[purge, dry\_run] As in \code{POST imports/\{id\}/remove}.
\end{params}
\raises{\code{400} no source.}

\apientry{GET demo/removal}{What removing the demo would remove}
\begin{lstlisting}[style=signature]
GET /api/v1/demo/removal
\end{lstlisting}
A dry run of deleting the demo's cohort \code{DEMO}, which removes the demo with it
(\cref{chap:quickstart}). Admins.
\returns{the counts of what would be removed.}

\apientry{GET mappings}{The saved column mappings}
\begin{lstlisting}[style=signature]
GET /api/v1/mappings
\end{lstlisting}
The column mappings saved for imports. Admins.
\returns{\code{\{"items": [\{name, kind, columns\}]\}}.}

\apientry{POST mappings}{Save a column mapping}
\begin{lstlisting}[style=signature]
POST /api/v1/mappings   {"name": NAME, "kind": "table", "columns": {FIELD: HEADER, ...}}
\end{lstlisting}
Saves a column mapping under a name, replacing one of the same name. Admins.
\begin{params}
\item[name] The mapping's name (not empty).
\item[kind] \code{table}.
\item[columns] Each catalog field and the header that holds it.
\end{params}
\returns{\code{\{\}}.}
\raises{\code{400} no name, or another kind.}

\apientry{DELETE mappings/\{name\}}{Delete a column mapping}
\begin{lstlisting}[style=signature]
DELETE /api/v1/mappings/{name}
\end{lstlisting}
Deletes a saved mapping (\code{name} is its name). Admins.
\returns{\code{\{\}}.}

\section{Connected databases}
\label{sec:http-api-connections}

\apientry{GET connections}{The connected databases}
\begin{lstlisting}[style=signature]
GET /api/v1/connections
\end{lstlisting}
The connected databases (\cref{chap:connecting}), with the persons linked; no secrets (only
the name of a password file). Admins.
\returns{\code{\{"items": [...]\}}.}

\apientry{POST connections}{Connect a database}
\begin{lstlisting}[style=signature]
POST /api/v1/connections   {"name": NAME, "kind": "omop" | "sql", "url": URL,
                            "password": TEXT, "settings": {...}}
\end{lstlisting}
Connects a database, as \cmd{medorama connections add} does, and checks that it answers.
Admins. Audited (\code{connect}).
\begin{params}
\item[name] Letters, digits, \code{-} and \code{\_}, starting with a letter (criteria write it:
  \code{name.lactate}); not the name of an imported database.
\item[kind] \code{omop} (an OMOP CDM database) or \code{sql} (views written in SQL).
\item[url] The database's address, without its password. A SQLite file
  (\code{sqlite:PATH}) must lie inside \code{data\_roots} and must not be the catalog itself.
\item[password] Optional: kept in \file{<home>/secrets/<name>.password} (mode 0600), never in
  the catalog, never sent back.
\item[settings] \code{schema}, \code{vocabulary\_schema}, \code{timezone} (of the database's
  times; when none is given, the catalog's time zone is written into the connection as it is
  added, and stays fixed from then on), \code{identifier} (what OMOP's
  \code{person\_source\_value} holds: \code{mrn}, the default, \code{patient\_id} or
  \code{none}: each person then becomes a patient of its own, with the patient ID
  \code{<connection>:<person key>}), \code{mrn\_authority}, \code{timeout\_s} (0: 60),
  \code{max\_rows} (per criterion, and for the persons a sync reads; 0: 1\,000\,000); for kind
  \code{sql} the views \code{persons} and \code{entries} (required), \code{visits}, and
  \code{version}.
\end{params}
\returns{\code{\{id, name, views, version\}}.}
\raises{\code{400} an invalid name, kind or URL, a file outside \code{data\_roots}, missing
views, a view that names a missing table (for OMOP every table its views read must exist, even
empty: \code{cdm\_source} is in effect required), or a database that does not answer.}

\apientry{POST connections/\{id\}/check}{Whether a connected database answers}
\begin{lstlisting}[style=signature]
POST /api/v1/connections/{id}/check
\end{lstlisting}
Asks a connected database for each of its views. Admins.
\returns{\code{\{"ok": true, "views", "version", "ms"\}}, or
\code{\{"ok": false, "error": TEXT\}}.}

\apientry{POST connections/\{id\}/sync}{Link a connected database's persons}
\begin{lstlisting}[style=signature]
POST /api/v1/connections/{id}/sync
\end{lstlisting}
Links every person of a connected database to the catalog's patients, reading them in one
query within \code{max\_rows}. Admins. Audited (\code{link}) when persons were linked, held or
found without valid identifiers.
\returns{\code{\{persons, linked\}}.}
\raises{\code{400} a database with more persons than \code{max\_rows} (raise \code{max\_rows},
or let searches link the persons they find).}

\apientry{DELETE connections/\{id\}}{Disconnect a database}
\begin{lstlisting}[style=signature]
DELETE /api/v1/connections/{id}
\end{lstlisting}
Disconnects a database: its links and what it added go, and the password file the page made
for it. Admins. Audited (\code{disconnect}).
\returns{\code{\{\}}.}

\section{Cohorts}
\label{sec:http-api-cohorts}

A cohort's criteria are a JSON object \code{\{"version": 1, "include": EXPR, "exclude":
[EXPR...]\}}, where an expression is a criterion (\code{\{"kind": ...\}}) or a group
(\code{\{"op": "and" | "or" | "not", "of": [EXPR...]\}}); the kinds of criteria and their fields
are those of \cref{chap:cohorts}. A cohort is seen and changed as it is shared: its owner may do
everything, a user it is shared with may read it, and edit it when the share allows.

\apientry{POST cohorts/search}{The patients that meet criteria}
\begin{lstlisting}[style=signature]
POST /api/v1/cohorts/search   {"criteria": {...}, "source": {...}, "within_source_of": ID,
                               "limit": N, "offset": N}
\end{lstlisting}
The patients that meet the criteria, with why. Any user.
\begin{params}
\item[criteria] The criteria.
\item[source] Only among a source's patients: \code{\{"op": "within" | "union" |
  "intersection" | "difference", "of": [COHORT...]\}}.
\item[within\_source\_of] Or among those of a stored cohort's source (its criteria being
  edited).
\item[limit, offset] The page (default 100 from 0).
\end{params}
\returns{\code{\{total, items, ids, excluded\_total, excluded\}}: \code{ids} every patient
found; \code{excluded} up to 50 of the patients the exclusions took out.}

\apientry{POST cohorts/parse}{Criteria written as text}
\begin{lstlisting}[style=signature]
POST /api/v1/cohorts/parse   {"text": TEXT}
\end{lstlisting}
Reads criteria written as text (\cref{chap:cohorts}). Any user.
\returns{\code{\{"ok": true, "expr": EXPR, "words": TEXT\}}, or
\code{\{"ok": false, "error": TEXT\}}.}
\begin{lstlisting}[style=shell]
curl -s -b jar -H 'X-Medorama: 1' -H 'Content-Type: application/json' \
  -d '{"text": "lactate > 4 and has icu"}' http://localhost:8090/api/v1/cohorts/parse
\end{lstlisting}

\apientry{GET cohorts}{The cohorts}
\begin{lstlisting}[style=signature]
GET /api/v1/cohorts
\end{lstlisting}
The cohorts the user may see, each with its \code{definition} in words. Any user.
\returns{\code{\{"items": [...]\}}.}

\apientry{POST cohorts}{Make a cohort}
\begin{lstlisting}[style=signature]
POST /api/v1/cohorts   {"name": NAME, "description": TEXT, "criteria": {...},
                        "source": {...}, "patients": [ID...]}
\end{lstlisting}
Makes a cohort, owned by the user. Any user. Audited (\code{cohort\_create}).
\begin{params}
\item[name] Its name (required).
\item[description] Optional.
\item[criteria] Its definition.
\item[source] A subcohort's parent or the cohorts it combines, as in \code{POST cohorts/search};
  none: the whole catalog.
\item[patients] The patients to put in now; none: all those the definition selects.
\end{params}
\returns{the cohort, as \code{GET cohorts/\{id\}}.}

\apientry{GET cohorts/\{id\}}{A cohort}
\begin{lstlisting}[style=signature]
GET /api/v1/cohorts/{id}
\end{lstlisting}
A cohort, with its definition and its shares. Any user it is shared with.
\raises{\code{404} no such cohort, or one the user may not see.}

\apientry{PATCH cohorts/\{id\}}{Change a cohort}
\begin{lstlisting}[style=signature]
PATCH /api/v1/cohorts/{id}   {"name", "description", "criteria", "add", "add_search",
                              "remove", "shares"}
\end{lstlisting}
Changes a cohort. Its editors; the shares, its owner only. Every field is optional. Audited
(\code{cohort\_change}, with what changed).
\begin{params}
\item[name, description, criteria] The new values.
\item[add] Patients to add: those the definition selects count as found by it, the others as
  added manually.
\item[add\_search] And every patient this search finds: \code{\{q, has, cohort\}}, as in
  \code{GET patients}.
\item[remove] Members to take out: those the definition selected stay out when it is refreshed.
\item[shares] Who else may see it, replacing the list:
  \code{[\{"user\_id": ID, "can\_edit": BOOL\}]}.
\end{params}
\returns{the cohort.}

\apientry{DELETE cohorts/\{id\}}{Delete a cohort}
\begin{lstlisting}[style=signature]
DELETE /api/v1/cohorts/{id}
\end{lstlisting}
Deletes a cohort. Its owner. Audited (\code{cohort\_delete}). The demo's cohort \code{DEMO} is
deleted by an admin only, and the demo goes with it (audited \code{demo\_remove}).
\returns{\code{\{\}}; for the demo's cohort \code{\{"removed": ...\}}, what was removed.}
\raises{\code{403} the demo's cohort by a user who is not an admin.}

\apientry{POST cohorts/\{id\}/clone}{Copy a cohort}
\begin{lstlisting}[style=signature]
POST /api/v1/cohorts/{id}/clone   {"name": NAME}
\end{lstlisting}
A copy of a cohort the user may see, owned by the user. Any user it is shared with.
\returns{the new cohort.}

\apientry{GET cohorts/\{id\}/history}{A cohort's changes}
\begin{lstlisting}[style=signature]
GET /api/v1/cohorts/{id}/history
\end{lstlisting}
A cohort's changes, latest first. Any user it is shared with.
\returns{\code{\{"items": [...]\}}.}

\apientry{GET cohorts/\{id\}/members}{A cohort's members}
\begin{lstlisting}[style=signature]
GET /api/v1/cohorts/{id}/members?limit=&offset=
\end{lstlisting}
A cohort's members. Any user it is shared with.
\begin{params}
\item[limit, offset] The page (default 100 from 0).
\end{params}
\returns{\code{\{total, items\}}.}

\apientry{POST cohorts/\{id\}/refresh}{Apply a cohort's definition again}
\begin{lstlisting}[style=signature]
POST /api/v1/cohorts/{id}/refresh   {"dry_run": BOOL, "join": [ID...], "leave": BOOL}
\end{lstlisting}
What the cohort's definition selects now, applied unless a dry run. Its editors. Audited
(\code{cohort\_refresh}) unless a dry run.
\begin{params}
\item[dry\_run] \code{true}: list who would join and who would leave (up to 500 of each),
  change nothing.
\item[join] Add only these of the patients who would join; none: all of them.
\item[leave] Whether the members the definition no longer selects leave (default \code{true});
  \code{false}: only new patients are added.
\end{params}
\returns{\code{added} and \code{removed}, and for a dry run the patients who would join and
leave.}

\apientry{GET cohorts/\{id\}/export}{Export a cohort}
\begin{lstlisting}[style=signature]
GET /api/v1/cohorts/{id}/export?what=&privacy=
\end{lstlisting}
A cohort as a file. Any user it is shared with. Through the gate in Privacy Mode, or with
\code{privacy=1} (\code{deid=1}, its earlier name).
\begin{params}
\item[what] \code{members} (the default: the members, CSV), \code{criteria} (the definition,
  JSON), \code{data} (the members' data, a long CSV) or \code{xlsx} (all of it, a sheet per
  database).
\item[privacy] \code{1} or \code{true}: through the privacy gate.
\end{params}
\returns{an attachment.}
\begin{lstlisting}[style=shell]
curl -s -b jar -o members.csv 'http://localhost:8090/api/v1/cohorts/2/export?privacy=1'
\end{lstlisting}

\apientry{GET cohorts/people}{Whom a cohort can be shared with}
\begin{lstlisting}[style=signature]
GET /api/v1/cohorts/people
\end{lstlisting}
The other enabled users, names only. Any user.
\returns{\code{\{"items": [\{id, username, display\_name\}]\}}.}

\section{Actions and jobs}
\label{sec:http-api-work}

\apientry{GET actions}{The actions}
\begin{lstlisting}[style=signature]
GET /api/v1/actions
\end{lstlisting}
The actions that can be run on recordings (\cref{chap:actions}), and the problems found reading
their descriptions. Any user.
\returns{\code{\{actions, problems\}}.}

\apientry{POST actions/reload}{Read the actions again}
\begin{lstlisting}[style=signature]
POST /api/v1/actions/reload
\end{lstlisting}
Reads the actions' descriptions again from their folders. Admins.
\returns{\code{\{actions, problems\}}.}

\apientry{POST runs}{Run actions}
\begin{lstlisting}[style=signature]
POST /api/v1/runs   {"actions": [NAME...], "selection": {...}, "params": {...},
                     "force": BOOL, "priority": N}
\end{lstlisting}
Queues a run: the actions on the recordings selected. Analysts and admins. Audited
(\code{run}).
\begin{params}
\item[actions] The actions' names.
\item[selection] Which recordings: \code{\{"encounters": [ID...]\}},
  \code{\{"recordings": [ID...]\}}, \code{\{"search": \{...\}\}} (a search of
  \code{GET encounters}), \code{\{"cohort": ID\}} (the encounters of its members) or
  \code{\{"all": true\}}.
\item[params] The actions' parameters.
\item[force] \code{true}: run again where already done on the unchanged recording.
\item[priority] 0 to 100 (default 0).
\end{params}
\returns{\code{\{run, recordings, queued, done\_already, pending\_already, not\_applicable\}}.}
\raises{\code{403} without the analyst or admin role.}

\apientry{GET jobs}{The jobs}
\begin{lstlisting}[style=signature]
GET /api/v1/jobs?state=&run=&limit=
\end{lstlisting}
The jobs, and their counts per state. Any user. A job's result (its summary) is not among
them: it comes with the encounter's recordings (\code{GET encounters/\{id\}}), and with
\code{GET recordings/\{id\}}.
\begin{params}
\item[state] Only the jobs in this state (\code{queued}, \code{running}, \code{done},
  \code{failed}, \code{cancelled}).
\item[run] Only the jobs of this run.
\item[limit] 1 to 2000 (default 200).
\end{params}
\returns{\code{\{items, counts\}}.}

\apientry{POST jobs/cancel}{Cancel jobs}
\begin{lstlisting}[style=signature]
POST /api/v1/jobs/cancel   {"job": ID} | {"run": ID} | {}
\end{lstlisting}
Cancels the queued and running jobs named: one job, the jobs of a run, or with \code{\{\}} all
of them (admins only). Analysts cancel only the jobs of their own runs, admins any. Analysts and
admins.
\returns{\code{\{"cancelled": N\}}.}
\raises{\code{403} a viewer, or \code{\{\}} by a user who is not an admin.}

\apientry{GET history}{One's history}
\begin{lstlisting}[style=signature]
GET /api/v1/history
\end{lstlisting}
The encounters the user opened in the encounter panel (each opening a
\code{GET encounters/\{id\}}), most recently opened first, with the number of views (up to 500).
Patient pages are not recorded. Only the user sees it. Any user.
\returns{\code{\{"items": [...]\}}.}

\apientry{DELETE history/\{id\}}{Forget one encounter of one's history}
\begin{lstlisting}[style=signature]
DELETE /api/v1/history/{id}
\end{lstlisting}
Removes an encounter (\code{id}, its number) from the user's history. Any user.
\returns{\code{\{"removed": N\}}.}

\apientry{DELETE history}{Forget one's history}
\begin{lstlisting}[style=signature]
DELETE /api/v1/history
\end{lstlisting}
Clears the user's history. Any user.
\returns{\code{\{"removed": N\}}.}

\section{Users and the audit log}
\label{sec:http-api-users}

\apientry{GET users}{The users}
\begin{lstlisting}[style=signature]
GET /api/v1/users
\end{lstlisting}
The users (\cref{chap:users-privacy}). Admins.
\returns{\code{\{"items": [...]\}}.}

\apientry{POST users}{Add a user}
\begin{lstlisting}[style=signature]
POST /api/v1/users   {"username": NAME, "display_name": TEXT, "role": ROLE, "phi": BOOL,
                      "password": PASSWORD}
\end{lstlisting}
Adds a user. Admins. Audited (\code{user\_added}). With proxy sign-in, the password may be left
out: the account then gets a random password that nobody knows, since the institution's sign-in
names the user (\cref{sec:sign-in-accounts}).
\begin{params}
\item[username] The user name.
\item[display\_name] Optional.
\item[role] \code{viewer}, \code{analyst} or \code{admin}.
\item[phi] The PHI permission (default \code{false}).
\item[password] At least 10 characters. Required with password sign-in; optional with proxy
  sign-in.
\end{params}
\returns{\code{\{"id": ID\}}.}
\raises{\code{400} an unknown role, a name taken, a password too short.}
\begin{lstlisting}[style=shell]
curl -s -b jar -H 'X-Medorama: 1' -H 'Content-Type: application/json' \
  -d '{"username": "reader1", "role": "viewer", "password": "a-long-passphrase"}' \
  http://localhost:8090/api/v1/users
\end{lstlisting}

\apientry{PATCH users/\{id\}}{Change a user}
\begin{lstlisting}[style=signature]
PATCH /api/v1/users/{id}   {"display_name", "role", "phi", "disabled", "password", "sign_out"}
\end{lstlisting}
Changes a user. Every field is optional. Admins. Audited (\code{user\_changed}, with what
changed).
\begin{params}
\item[display\_name, role, phi, disabled] The new values.
\item[password] A new password (a reset).
\item[sign\_out] \code{true}: end the user's sessions on every device.
\end{params}
\returns{\code{\{\}}.}
\raises{\code{400} an unknown role, or a change that would leave the catalog without an enabled
admin.}

\apientry{DELETE users/\{id\}}{Delete a user}
\begin{lstlisting}[style=signature]
DELETE /api/v1/users/{id}
\end{lstlisting}
Deletes a user. Admins. Audited (\code{user\_deleted}).
\returns{\code{\{\}}.}
\raises{\code{400} the last enabled admin.}

\apientry{GET audit}{The audit log}
\begin{lstlisting}[style=signature]
GET /api/v1/audit?limit=
\end{lstlisting}
The latest entries of the audit log, latest first (\cref{chap:users-privacy}). Admins. In
Privacy Mode an entry's \code{target} and \code{detail} are \code{[Withheld]} (a search's text
may name a patient).
\begin{params}
\item[limit] 1 to 10\,000 (default 500).
\end{params}
\returns{\code{\{"items": [\{time, username, event, target, detail\}]\}}.}

\apientry{GET audit/verify}{Whether the audit log's seals hold}
\begin{lstlisting}[style=signature]
GET /api/v1/audit/verify
\end{lstlisting}
Computes the chain of every seal of the audit log again, counts the entries and seals the
database numbered against those left, and compares the seals with the export file
(\code{audit\_export}), when it is set: each exported seal must be in the database, with the
same hash. Admins.
\returns{\code{\{seals, entries, unsealed, broken, exported\_seals\}}: \code{broken}
\code{null}, or the first seal whose chain differs, what is missing, or an exported seal not in
the database, and how; \code{exported\_seals} the seals of the export file (\code{null}
without one).}
\begin{lstlisting}[language=JSON]
{"seals": 14, "entries": 1830, "unsealed": 12, "broken": null, "exported_seals": 14}
\end{lstlisting}

\section{The waveform viewer}
\label{sec:http-api-viewer}

\apientry{GET viewer}{The viewer's status}
\begin{lstlisting}[style=signature]
GET /api/v1/viewer
\end{lstlisting}
Whether a VitalStream viewer is available under the catalog (\cref{chap:viewers-in-medorama}).
Any user.
\returns{\code{\{mode, flavor, running, reason, url\}}: \code{mode} \code{embedded} (the
catalog runs a viewer, at \code{url} \code{viewer/}), \code{link} (\code{viewer\_url} is set:
\code{url} is it) or \code{none} (\code{reason} says why); \code{flavor} the viewer run
(\code{cpp}, \code{rs}, \code{py}, or \code{custom} for a program named in the
settings).}
\begin{lstlisting}[language=JSON]
{"mode": "embedded", "flavor": "rs", "running": true, "reason": "", "url": "viewer/"}
\end{lstlisting}

\apientry{POST viewer/start}{Start the viewer now}
\begin{lstlisting}[style=signature]
POST /api/v1/viewer/start
\end{lstlisting}
Starts the embedded viewer now, so that the first view is quick. Any user.
\returns{the status, as \code{GET viewer}.}
\raises{\code{404} no viewer runs under this catalog (\code{reason}); \code{503} it could not
be started.}
